\documentclass{article}

\usepackage{./styles/main}
\graphicspath{{./figure/}}
\addbibresource{src/references.bib}

\begin{document}

% ===================== ABSTRACTS FIRST =====================
% TODO: Replace placeholders with your real topic, goals, results, metrics.
\begin{enabstract}
  This thesis develops and evaluates a hybrid approach to handwritten text image classification.
  The goal is to improve accuracy on small and medium datasets by expanding the feature space with
  outputs of a Probabilistic Neural Network (PNN), followed by training standard machine-learning
  classifiers.

  The contribution is a reproducible pipeline that couples a PNN feature generator with linear
  and ensemble classifiers, reducing hyperparameter search while improving robustness to handwriting
  variation. Experiments on a real biomedical dataset show a $5$--$8\,\%$ F1 improvement over a
  baseline.

  \textbf{Total length}: 70 pages, 33 figures, 34 references.

  \textbf{Keywords}: probabilistic neural network, classification, handwriting, feature expansion, machine learning.
\end{enabstract}

\begin{uaabstract}
  Бакалаврська кваліфікаційна робота присвячена розробленню та експериментальному
  дослідженню гібридного підходу до класифікації зображень рукописного тексту. Мета роботи —
  підвищити точність класифікації у випадку малих і середніх навчальних вибірок за рахунок
  розширення простору ознак виходами ймовірнісної нейронної мережі та подальшого навчання
  класичних алгоритмів машинного навчання.

  Наукова новизна полягає у поєднанні PNN як генератора інформативних ознак із
  лінійними та ансамблевими класифікаторами, що дає змогу зменшити розмірність пошуку
  гіперпараметрів і підвищити стійкість до варіативності почерку. Практична цінність — у створенні
  відтворюваного конвеєра, який дозволяє порівнювати моделі, підбирати гіперпараметри та
  автоматично формувати звіт.

  Експерименти на реальному біомедичному наборі демонструють підвищення F1-міри на
  $5$–$8\,\%$ порівняно з базовою лінією.

  \textbf{Загальний обсяг}: 70 сторінок, 33 рисунки, 34 посилання.

  \textbf{Ключові слова}: ймовірнісна нейронна мережа, класифікація, рукописний текст, розширення
  вибірки, машинне навчання.
\end{uaabstract}

% ===================== TABLE OF CONTENTS =====================
\clearpage
\tableofcontents
\clearpage

% ===================== ABBREVIATIONS =====================
% TODO: Adjust to your topic; keep only what you use in the text.
\section*{LIST OF ABBREVIATIONS}
\addcontentsline{toc}{section}{LIST OF ABBREVIATIONS}
\begin{description}
  \item[API] Application Programming Interface.
  \item[AI] Artificial Intelligence.
  \item[ML] Machine Learning.
  \item[DL] Deep Learning.
  \item[PNN] Probabilistic Neural Network.
  \item[UMAP] Uniform Manifold Approximation and Projection.
  \item[t-SNE] t-distributed Stochastic Neighbor Embedding.
  \item[CART] Classification and Regression Tree.
\end{description}
\clearpage

% ===================== INTRODUCTION =====================
\section*{INTRODUCTION}
\addcontentsline{toc}{section}{INTRODUCTION}
% (Why, What, How, Results, Structure)
The digitization of documents and manuscripts creates demand for accurate and reproducible methods
for classifying textual images. The \textbf{goal of the work} is to improve classification accuracy under
data constraints.

\textbf{Research tasks:}
\begin{itemize}
  \item review and systematize methods of segmentation and feature extraction \cite{otsu1979,canny1986};
  \item justify a hybrid approach with feature expansion via a PNN;
  \item develop a training, validation, and evaluation pipeline;
  \item conduct experiments on real data and analyze the results;
  \item formulate conclusions, limitations, and directions for future research.
\end{itemize}

The paper \cite{lecun1998gradient} classically formulates the ``gradient-based learning'' approach to document
recognition; convolutional networks \cite{krizhevsky2012imagenet,he2016deep} and optimization methods \cite{kingma2014adam}
form the technical foundation of modern solutions. For visualization of high-dimensional features,
t-SNE \cite{tsne2008} and UMAP \cite{umap2018} are commonly used.

\textbf{Structure of the thesis.} Chapter~1 provides the overview and problem statement. Chapter~2
describes the mathematical model and algorithmic components. Chapter~3 covers implementation and
experiments. The work concludes with general findings.

% ===================== CHAPTER 1 =====================
\section{GENERAL PROVISIONS}

\subsection{Problem context and problem statement}
The objective is to build a classifier $f:\mathcal{X}\to\mathcal{Y}$, where $\mathcal{X}$ is the space of
character/word images and $\mathcal{Y}$ is the set of classes. \textit{Input} data: grayscale images,
non-uniform scanning conditions, and handwriting variability. \textit{Output}: a class label and,
when available, confidences.

\subsection{Domain overview}
Contour-based methods and thresholding remain fundamental tools for preprocessing
\cite{otsu1979,canny1986}. Deep models \cite{lecun1998gradient,krizhevsky2012imagenet,he2016deep}
provide high quality but require large datasets and careful tuning.

\subsubsection{Basic processing stages}
\begin{enumerate}
  \item normalization and brightness leveling;
  \item binarization (e.g., Otsu \cite{otsu1979});
  \item edge detection (Canny \cite{canny1986});
  \item character/word segmentation; feature extraction.
\end{enumerate}

\paragraph{Illustrative example.}
A schematic of the pipeline is shown in Fig.~\ref{fig:pipeline}.

\begin{figure}[htbp]
  \centering
  \includegraphics[width=0.9\linewidth]{example-figure.jpg}
  \caption{Schematic of the proposed handwritten text classification pipeline}
  \label{fig:pipeline}
\end{figure}

\subsection{Approach formulation}
Hypothesis: \emph{PNN outputs as generalized posterior estimates} can serve as a compact and
stable representation for classical classifiers. Thus, we seek a mapping
$g:\mathcal{X}\to\mathbb{R}^k$ (PNN features) and a classifier $h:\mathbb{R}^k\to\mathcal{Y}$ such that
$f(x)=h(g(x))$.

% ===================== CHAPTER 2 =====================
\section{INFORMATIONAL AND MATHEMATICAL FOUNDATIONS}

\subsection{Preprocessing and feature extraction}
Let $I:\Omega\to[0,1]$ denote a normalized image. Adaptive thresholding
\[
  T^\ast=\arg\max_T \ \mathrm{BetweenClassVar}(I,T)
\]
implements Otsu's criterion \cite{otsu1979}.
Gradient operators are used for boundary detection \cite{canny1986}.

\subsection{Hybrid model: PNN\,$\to$\,classifier}
Let $\phi_\sigma(x)$ be the vector of PNN outputs with smoothing parameter $\sigma$. We then
train $h$ (SVM, Random Forest, or logistic regression) on $\phi_\sigma(x)$.
For parameter optimization we use $k$-fold cross-validation.
\begin{equation}
  \label{eq:objective}
  \hat{\theta}=\arg\max_{\theta} \ \frac{1}{k}\sum_{i=1}^{k}\mathrm{F1}\bigl(h_{\theta}(\phi_\sigma(X_{-i})),\,Y_{-i}\bigr).
\end{equation}
Equation~\eqref{eq:objective} minimizes the \emph{generalization error} and protects against
overfitting.

\paragraph{Quality metrics.}
We use Accuracy, Precision, Recall, and F1. An example layout is shown in Table~\ref{tab:metrics}.
\begin{table}[htbp]
  \caption{Comparison of quality metrics for three approaches}
  \label{tab:metrics}
  \centering
  \begin{tabular}{|l|c|c|c|}
    \hline
    Method   & Precision & Recall & F1   \\
    \hline
    Baseline & 0.84      & 0.81   & 0.82 \\
    \hline
    PNN+SVM  & 0.89      & 0.87   & 0.88 \\
    \hline
    PNN+RF   & 0.90      & 0.88   & 0.89 \\
    \hline
  \end{tabular}
\end{table}

\subsection{Visualization and interpretation}
To inspect cluster structure, we apply t-SNE \cite{tsne2008} and UMAP \cite{umap2018}.
\begin{enumerate}
  \item t-SNE reveals local structure well but is sensitive to hyperparameters;
  \item UMAP more often preserves global structure and scales better.
\end{enumerate}

\subsection{General notes on reproducibility}
\begin{itemize}
  \item fix random seeds;
  \item version datasets and the execution environment;
  \item log hyperparameters and metrics in a unified format (CSV/JSON).
\end{itemize}

% ===================== CHAPTER 3 =====================
\section{SOFTWARE AND TECHNICAL SUPPORT}

\subsection{System architecture}
The system consists of modules: \textit{io} (data loading), \textit{preprocess},
\textit{features} (PNN), \textit{models} (classifiers), and \textit{eval} (metrics, reports).
% TODO: add a class/module diagram of your implementation to the report.

\subsection{Runtime environment and requirements}
Python~3.x; libraries for image processing and ML; sufficient RAM for training.
For optimization we use Adam \cite{kingma2014adam} and cyclic learning-rate schedules as needed.

\subsection{Experimental protocol}
\begin{itemize}
  \item split into train/validation/test with no overlap of samples from the same author;
  \item stratified sampling for class balance;
  \item $k$-fold cross-validation to tune $\sigma$ (PNN) and the hyperparameters of $h$;
  \item statistical significance testing (e.g., paired t-test on F1).
\end{itemize}

\subsection{Threats to validity}
Labeling errors, data drift, and test leakage through preprocessing. It is necessary to
check sensitivity to parameters and generalization to other corpora.

% ===================== GENERAL CONCLUSIONS =====================
\section*{GENERAL CONCLUSIONS}
\addcontentsline{toc}{section}{GENERAL CONCLUSIONS}
A hybrid PNN\,$\to$\,classifier approach is proposed that increases F1 in data-limited scenarios.
A reproducible experimentation pipeline was developed; a comparative analysis was performed; and
practical recommendations for model and hyperparameter selection were formulated.

% ===================== BIBLIOGRAPHY =====================
\clearpage
\printbibliography[heading=bibintoc,title={REFERENCES}]

% ===================== APPENDICES =====================
\appendix

\section*{Appendix A}
\addcontentsline{toc}{section}{Appendix A}
% TODO: Move additional tables/illustrations and data examples here.

\section*{Appendix B}
\addcontentsline{toc}{section}{Appendix B}
% TODO: Link to the repository; reproduction instructions (CLI, configuration, versions).

\end{document}
